% GDZ PPN236018647,native591lower;HKUST592.
\providecommand{\DiaryMoonRight}{\tikz[baseline=0pt,x=1ex,y=1ex]{\draw[line width=.4pt](0,1.3)..controls(1.16,1.3)and(1.16,0)..(0,0)..controls(.8,.08)and(.8,1.22)..(0,1.3);}}
\providecommand{\DiaryCeres}{\tikz[baseline=0pt,x=1ex,y=1ex]{\draw[line width=.45pt](.05,.65)..controls(.45,.2)and(.95,.65)..(.75,1.1)..controls(.58,1.53)and(.16,1.5)..(.02,1.22);\draw[line width=.45pt](.05,.65)--(.05,0);\draw[line width=.45pt](-.22,.3)--(.38,.3);}}
\NumberedItem{120}
\noindent\hspace*{\parindent}Theoriam motus Lunae aggressi sumus.
\par\noindent\hfill [1801] Aug.
\par\medskip\begingroup\small
\noindent\hspace*{\parindent}Die um diese Zeit entstandene \textit{Theorie der Bewegung des Mondes} ist im Nachlaß vorhanden und Werke VII, 1906, S. 611—639 abgedruckt. In einem Briefe an \textsc{Schumacher} vom 23. Januar 1842 (\textit{Briefwechsel zwischen Gauss und Schumacher} IV, Altona 1862, S. 51) schreibt \textsc{Gauss}:
\par\endgroup
\par\noindent\hspace*{\parindent}Eben im Sommer 1801 hatte ich mir vorgesetzt, ähnliche Arbeit über den \DiaryMoonRight\ [Mond] auszuführen, aber kaum hatte ich die theoretischen Vorarbeiten angefangen (denn diese sind es, auf welche in der Vorrede meiner \textit{Th[eoria] M[otus] C[orporum] C[oelestium]} angespielt wird), als das Bekanntwerden von \textsc{Piazzis} \DiaryCeres\ [Ceres] Beob[achtungen] mich in eine ganz andere Richtung zog.
\par\medskip\begingroup\small
\noindent\hspace*{\parindent}Die in diesem Briefe erwähnte Stelle aus der Vorrede zur \textit{Theoria Motus} ist oben bei der Nr. 119 abgedruckt. Augenscheinlich sind die Tagebucheintragungen Nr. 119 und Nr. 120 gleichzeitig gemacht.
\par\noindent\hfill\textsc{Brendel.}
\par\endgroup
